\chapter{Potassium (Blood)}

\section*{What Is This Test?}
The potassium blood test measures the concentration of potassium ions (K\textsuperscript{+}) in the blood serum or plasma. Potassium is the primary intracellular cation, with approximately 98\% of total body potassium located inside cells, and only 2\% in the extracellular fluid. The normal serum potassium concentration is maintained within a narrow range of 3.5--5.0 mEq/L (mmol/L). This tight regulation is essential because potassium is critical for maintaining the resting membrane potential of all excitable cells, including cardiac myocytes, skeletal muscle cells, and neurons. The sodium-potassium ATPase pump (Na\textsuperscript{+}/K\textsuperscript{+}-ATPase) actively transports potassium into cells against its concentration gradient, consuming a significant portion of cellular energy. Factors affecting potassium distribution include: plasma potassium concentration, insulin (promotes cellular uptake via Na\textsuperscript{+}/K\textsuperscript{+}-ATPase stimulation), catecholamines (beta-2 adrenergic agonists promote uptake; beta-blockers shift potassium out), acid-base status (acidosis shifts potassium out of cells; alkalosis shifts potassium into cells), plasma osmolality, and thyroid hormone. Renal excretion through the collecting duct accounts for 90--95\% of daily potassium elimination, with aldosterone being the major hormonal regulator. Serum potassium represents only about 0.4\% of total body potassium stores, so small shifts between compartments can cause dramatic changes in measured levels.

\section*{Alternative Names}
K; Potassium, Serum; Serum Potassium; Plasma Potassium; K+ Level; Kalium

\section*{Principle}
Serum potassium is measured using ion-selective electrode (ISE) potentiometry with a valinomycin-based membrane electrode. Valinomycin is a cyclic ionophore antibiotic that encapsulates potassium ions with high selectivity (selectivity coefficient K\textsuperscript{+}/Na\textsuperscript{+} of approximately 10,000:1), allowing them to traverse the hydrophobic membrane while excluding smaller ions like sodium and lithium. The potential difference across the membrane is proportional to the logarithm of potassium activity according to the Nernst equation. Both direct ISE (on blood gas analyzers, measuring potassium activity in undiluted plasma water) and indirect ISE (on automated chemistry analyzers after dilution) are used clinically. Direct ISE is not affected by pseudohyperkalemia from extreme thrombocytosis or leukocytosis. The measurement can also be performed by enzymatic methods based on potassium-dependent pyruvate kinase activity, which are used on some dry chemistry platforms (e.g., Vitros). Results are reported in mEq/L, which is numerically equivalent to mmol/L.

\section*{History}
The critical role of potassium in physiology was first recognized in the late 19th century, with Sydney Ringer (1883) demonstrating that potassium was essential for cardiac contraction in his experiments on frog hearts. The flame photometry method for potassium measurement was introduced in the 1940s-1950s and remained the standard for decades. The valinomycin-based potassium ion-selective electrode was developed in the 1970s by Simon and colleagues at ETH Z\"{u}rich, revolutionizing clinical measurement by eliminating the need for flame photometers. The discovery of the sodium-potassium ATPase pump by Jens Christian Skou in 1957 earned him the 1997 Nobel Prize in Chemistry. Wills and colleagues established modern reference ranges in the 1970s-1980s.

\section*{Why This Test Is Ordered}
\begin{itemize}
  \item Evaluation of cardiac arrhythmias (especially atrial and ventricular ectopy, prolonged QT, U waves, widened QRS, sine wave pattern), muscle weakness, paralysis, or paresthesias
  \item Monitoring of patients on medications that profoundly affect potassium homeostasis: diuretics (loop diuretics---furosemide, bumetanide---cause hypokalemia; potassium-sparing diuretics---spironolactone, eplerenone, amiloride, triamterene---cause hyperkalemia), ACE inhibitors, angiotensin receptor blockers (ARBs), digoxin (toxicity worsened by hypokalemia), succinylcholine, beta-blockers, heparin, NSAIDs, trimethoprim, pentamidine, calcineurin inhibitors (tacrolimus, cyclosporine)
  \item Monitoring during IV potassium replacement therapy to avoid overshoot hyperkalemia
  \item Evaluation of hypertension (hyperaldosteronism screening), acid-base disorders (metabolic acidosis elevates potassium; metabolic alkalosis lowers it), and renal function
  \item Acute kidney injury (AKI) and chronic kidney disease (CKD) monitoring---hyperkalemia is a major cause of mortality in ESRD
  \item Preoperative evaluation and perioperative monitoring
  \item Workup of adrenal disorders (Addison's disease, primary hypoaldosteronism, congenital adrenal hyperplasia)
  \item Monitoring tumor lysis syndrome in patients receiving chemotherapy for hematologic malignancies
  \item Evaluation of rhabdomyolysis (massive potassium release from damaged muscle)
\end{itemize}

\section*{Is This a Primary or Secondary Test}
Potassium measurement is a primary screening test included in the basic metabolic panel (BMP). It is one of the most critical components of electrolyte monitoring because both hypokalemia and hyperkalemia can be immediately life-threatening due to cardiac arrhythmia risk. Serial potassium monitoring is standard of care for patients receiving diuretics, IV fluids, or with renal impairment.

\section*{How This Test Is Performed}
A healthcare professional will draw blood from a vein in your arm. A tourniquet is applied, the skin is cleaned, and a needle is inserted into a vein. Blood is collected in a serum separator tube (gold-top) or lithium heparin tube (green-top). EDTA tubes (purple-top) must NOT be used for potassium measurement because EDTA contains potassium salts that will contaminate the sample and cause falsely elevated results. The tourniquet should be released as soon as possible---prolonged application and repeated fist clenching during venipuncture can cause a 0.5--1.0 mEq/L increase in potassium from local tissue release. The sample should be centrifuged promptly (within 30 minutes)---prolonged contact between serum/plasma and red blood cells allows potassium to leak out of cells, causing pseudohyperkalemia. Results are available within 30 minutes to 1 hour. For stat determinations, whole blood potassium can be measured on a blood gas analyzer using direct ISE in under 5 minutes.

\section*{What Preparation Is Needed}
No fasting is required for a potassium level. However, inform your healthcare provider of all medications (especially diuretics, ACE inhibitors, ARBs, potassium supplements, NSAIDs, heparin, beta-blockers), as these directly affect potassium levels. Avoid fist clenching during phlebotomy (can increase potassium by 0.5--1.0 mEq/L). The sample must be delivered to the laboratory promptly---delays >2 hours at room temperature without centrifugation cause potassium leakage from red blood cells (pseudohyperkalemia). If a delay is unavoidable, the sample can be centrifuged and the serum/plasma separated within 30 minutes of collection. Refrigerated storage at 2--8\textdegree C before centrifugation is NOT recommended for potassium---cold temperature inhibits Na\textsuperscript{+}/K\textsuperscript{+}-ATPase, causing cellular potassium leakage.

\section*{Contraindications and Precautions}
No absolute contraindications to venipuncture. Critical pre-analytical considerations for potassium: (1) Hemolysis---the most common cause of pseudohyperkalemia, as red blood cells contain potassium at approximately 100 mEq/L (30 times serum concentration); hemolyzed samples must be rejected and redrawn. Even slight hemolysis (plasma/serum pink-tinged) can significantly elevate potassium. (2) Prolonged tourniquet use with fist clenching causes 0.5--1.0 mEq/L increase. (3) EDTA tube contamination (purple-top)---EDTA chelates divalent cations and contains potassium salt, causing grossly erroneous results. (4) Extreme thrombocytosis (platelets >1,000,000/$\mu$L)---potassium released from platelets during clotting process causes pseudohyperkalemia in serum samples; in these patients, plasma (lithium heparin) potassium should be measured. (5) Extreme leukocytosis (WBC >100,000/$\mu$L)---leukemic cells are fragile and may release potassium during processing. (6) Delayed separation---potassium leaks from red blood cells over time. (7) Recent exercise---transient hyperkalemia occurs from potassium release during muscle contraction; resolves within minutes of rest. (8) Hyperventilation during phlebotomy can cause respiratory alkalosis with intracellular potassium shift (pseudohypokalemia). (9) Specimens drawn from an arm receiving IV fluids with potassium (falsely elevated) or without potassium (dilutional effect). (10) Digoxin toxicity: hypokalemia potentiates digoxin binding to Na\textsuperscript{+}/K\textsuperscript{+}-ATPase, increasing cardiotoxicity; mild hypokalemia (3.0--3.5 mEq/L) in a patient on digoxin warrants immediate potassium repletion.

\section*{Risks}
Minimal risk from venipuncture. Mild pain or bruising at the puncture site is common and typically resolves within 1--2 days. Rare complications: hematoma, infection, nerve injury, vasovagal syncope. No significant risk of blood loss (3--5 mL blood drawn). If multiple attempts are required, there is a slightly higher chance of hemolysis from traumatic venipuncture, which would affect potassium accuracy.

\section*{What Happens After the Test}
A bandage is applied and should be kept for 15--30 minutes. Results are available within 30 minutes to 1 hour for routine testing, or within 5--15 minutes for stat/point-of-care testing. Critical potassium values requiring immediate clinical action: <2.5 mEq/L or >6.5 mEq/L. Potassium <2.5 mEq/L carries risk of respiratory muscle paralysis, rhabdomyolysis, and cardiac arrest (especially in patients with heart disease or on digoxin). Potassium >6.5 mEq/L carries risk of life-threatening cardiac arrhythmia (peaked T waves, loss of P waves, widened QRS progressing to sine wave pattern and ventricular fibrillation/asystole). ECG monitoring is warranted for potassium <3.0 mEq/L or >6.0 mEq/L. Emergency treatment of severe hyperkalemia includes IV calcium gluconate/chloride (cardioprotective, stabilizes myocardial membrane potential), IV insulin with dextrose (shifts potassium intracellularly), IV sodium bicarbonate (shifts potassium intracellularly in acidotic patients), beta-2 agonists (nebulized albuterol), and definitive removal with loop diuretics, sodium polystyrene sulfonate, patiromer, sodium zirconium cyclosilicate, or urgent hemodialysis.

\section*{Understanding Results}

\subsection*{Below Normal}
\begin{itemize}
  \item \textbf{Hypokalemia (K <3.5 mEq/L):} Present in approximately 20\% of hospitalized patients; up to 40\% of patients on thiazide diuretics.
  \item \textbf{Mild (3.0--3.4 mEq/L):} Often asymptomatic but may cause mild muscle weakness, fatigue, constipation. ECG may show T wave flattening and prominent U waves.
  \item \textbf{Moderate (2.5--2.9 mEq/L):} Muscle weakness, cramps, ileus, confusion, increased risk of cardiac arrhythmia (especially if digoxin present, if LVH, or if ischemic heart disease). ECG: ST depression, flattened T waves, prominent U waves, prolonged QT interval.
  \item \textbf{Severe (<2.5 mEq/L):} Profound weakness, rhabdomyolysis, respiratory failure (due to diaphragmatic and respiratory muscle paralysis), ileus, cardiac arrest from ventricular fibrillation or asystole.
  \item \textbf{Renal losses (most common cause):} Diuretic therapy (especially thiazide and loop diuretics); primary hyperaldosteronism (Conn's syndrome); secondary hyperaldosteronism (heart failure, cirrhosis, renal artery stenosis); Cushing's syndrome; Bartter syndrome, Gitelman syndrome; Liddle syndrome; renal tubular acidosis (type I and II); amphotericin B toxicity; cisplatin nephrotoxicity; hypomagnesemia (impaired renal potassium conservation---hypomagnesemia causes renal potassium wasting that is refractory to potassium replacement until magnesium is corrected); licorice ingestion (glycyrrhizic acid inhibits 11-$\beta$-HSD2, causing apparent mineralocorticoid excess).
  \item \textbf{Gastrointestinal losses:} Vomiting (metabolic alkalosis increases renal potassium excretion; gastric fluid itself contains only 5--10 mEq/L potassium---the hypokalemia is primarily from renal loss); diarrhea (colonic fluid contains 20--50 mEq/L potassium; secretory diarrhea from cholera, VIPoma, villous adenoma causes massive potassium depletion); nasogastric suction; laxative abuse; villous adenoma.
  \item \textbf{Transcellular shifts:} Insulin administration; beta-2 adrenergic agonists (albuterol, epinephrine); metabolic alkalosis; hypokalemic periodic paralysis (genetic sodium or calcium channel mutation); refeeding syndrome; treatment of megaloblastic anemia (rapid cellular uptake); barium toxicity; theophylline overdose.
  \item \textbf{Inadequate intake:} Rare as sole cause but contributory in alcoholism, anorexia nervosa, starvation, prolonged IV fluids without potassium supplementation; typically requires at least 10--14 days of near-zero intake before significant hypokalemia develops.
\end{itemize}

\subsection*{Above Normal}
\begin{itemize}
  \item \textbf{Hyperkalemia (K >5.0 mEq/L):} Present in 1--10\% of hospitalized patients; potentially the most immediately dangerous electrolyte abnormality due to cardiac conduction effects.
  \item \textbf{Mild (5.1--5.9 mEq/L):} May be asymptomatic. Ensure not pseudohyperkalemia from hemolysis, prolonged tourniquet, or delayed processing. ECG may be normal or show peaked T waves (tall, narrow, symmetric peaked T waves best seen in precordial leads V2--V4).
  \item \textbf{Moderate (6.0--6.4 mEq/L):} Risk of cardiac conduction delay. ECG: peaked T waves, prolonged PR interval, loss of P wave amplitude, widened QRS complex. Immediate intervention required: IV calcium, insulin + dextrose, and urgent nephrology consultation.
  \item \textbf{Severe (>6.5 mEq/L):} Life-threatening emergency. ECG progression: loss of P waves $\rightarrow$ widened QRS $\rightarrow$ sine wave pattern $\rightarrow$ ventricular fibrillation $\rightarrow$ asystole. Requires ICU admission, continuous ECG monitoring, and emergency dialysis if refractory to medical therapy.
  \item \textbf{Renal failure (most common cause):} Acute kidney injury; chronic kidney disease (eGFR <15 mL/min/1.73 m\textsuperscript{2}); ESRD on dialysis---hyperkalemia is the second most common cause of death in dialysis patients after cardiovascular disease. Type IV renal tubular acidosis (hyporeninemic hypoaldosteronism, common in diabetic nephropathy and tubulointerstitial disease).
  \item \textbf{Medications (very common):} ACE inhibitors, ARBs, potassium-sparing diuretics (spironolactone, eplerenone, amiloride, triamterene), NSAIDs (reduce renin and aldosterone), trimethoprim (blocks ENaC sodium channel like amiloride), calcineurin inhibitors (tacrolimus, cyclosporine---inhibit renal potassium secretion), heparin (inhibits aldosterone synthesis), succinylcholine (depolarizing muscle relaxant causing potassium efflux), digoxin toxicity (inhibits Na\textsuperscript{+}/K\textsuperscript{+}-ATPase), beta-blockers (reduce cellular potassium uptake).
  \item \textbf{Adrenal insufficiency (Addison's disease):} Primary adrenal insufficiency (autoimmune, tuberculosis, adrenal hemorrhage) causes aldosterone deficiency with hyponatremia, hyperkalemia, and metabolic acidosis.
  \item \textbf{Transcellular shifts:} Metabolic acidosis (inorganic/mineral acid acidosis---HCl, NH4Cl---causes greater shift than organic acidosis); insulin deficiency (diabetic ketoacidosis, unrecognized type 1 diabetes); hyperglycemic hyperosmolar state; rhabdomyolysis; tumor lysis syndrome (massive cellular destruction from chemotherapy in leukemia, lymphoma); massive blood transfusion (>10 units of stored blood---potassium leaks from stored RBCs); succinylcholine use; digitalis toxicity; beta-blocker overdose.
  \item \textbf{Excessive intake:} Oral or IV potassium supplements (especially in setting of reduced renal function); potassium-containing salt substitutes (KCl-based); large-volume transfusion; potassium penicillin G (1.7 mEq per million units).
  \item \textbf{Pseudohyperkalemia:} Hemolysis (most common---always check sample quality); fist clenching during phlebotomy; prolonged tourniquet; thrombocytosis (>1,000,000/$\mu$L); leukocytosis (>100,000/$\mu$L); delayed serum separation; EDTA tube contamination; familial pseudohyperkalemia (rare genetic red cell membrane defect with temperature-dependent potassium leak).
\end{itemize}

\section*{Normal Results}
\begin{itemize}
  \item \textbf{Adults:} 3.5--5.0 mEq/L (mmol/L)
  \item \textbf{Children (1--16 years):} 3.4--4.7 mEq/L
  \item \textbf{Infants (1--12 months):} 3.5--5.5 mEq/L
  \item \textbf{Newborns (0--7 days):} 3.7--5.9 mEq/L (slightly higher due to decreased renal function in first week of life)
  \item \textbf{Plasma (lithium heparin):} 3.5--5.0 mEq/L (approximately 0.1--0.4 mEq/L lower than serum due to potassium release from platelets during clotting)
  \item \textbf{Critical low value:} <2.5 mEq/L
  \item \textbf{Critical high value:} >6.5 mEq/L
\end{itemize}

\section*{Follow-up Tests}
\begin{itemize}
  \item \textbf{ECG:} Urgently in hyperkalemia >6.0 mEq/L or hypokalemia <3.0 mEq/L with cardiac symptoms. Look for peaked T waves, loss of P waves, widened QRS (hyperkalemia); flattened T waves, U waves, prolonged QT (hypokalemia).
  \item \textbf{Repeat potassium measurement:} To verify abnormal result (rule out pseudohyperkalemia from hemolysis). Use plasma (lithium heparin) tube if thrombocytosis is suspected.
  \item \textbf{Comprehensive metabolic panel (including BUN, creatinine, glucose):} Renal function assessment (hyperkalemia most commonly from renal failure); glucose check (hyperkalemia in DKA/hyperosmolar states).
  \item \textbf{Magnesium level:} Hypomagnesemia causes renal potassium wasting that is refractory to potassium replacement alone. Always check magnesium in refractory hypokalemia. Correct magnesium to >2.0 mg/dL before potassium replacement will be effective.
  \item \textbf{Arterial blood gas (ABG):} To evaluate acid-base status; pH change of 0.1 changes potassium by ~0.5 mEq/L in the opposite direction.
  \item \textbf{Plasma renin activity and serum aldosterone:} If hyperkalemia with hypertension is present (hyperaldosteronism screening); also in hypokalemia to investigate primary hyperaldosteronism (elevated aldosterone, suppressed renin).
  \item \textbf{Morning cortisol and ACTH stimulation test:} For suspected adrenal insufficiency causing hyperkalemia with hyponatremia.
  \item \textbf{24-hour urine potassium:} To differentiate renal from extrarenal potassium losses. Urine K >20 mEq/day in setting of hypokalemia indicates inappropriate renal potassium wasting.
  \item \textbf{Transtubular potassium gradient (TTKG):} <3 in hyperkalemia suggests appropriate renal response (or mineralocorticoid deficiency); >7 in hypokalemia suggests appropriate renal potassium conservation. Used less now due to incomplete understanding of collecting duct physiology.
  \item \textbf{Creatine kinase (CK):} If rhabdomyolysis is suspected as cause of hyperkalemia.
  \item \textbf{Uric acid, LDH, phosphate:} For tumor lysis syndrome monitoring in patients receiving chemotherapy.
\end{itemize}

\section*{References}
\begin{enumerate}
  \item Burtis CA, Ashwood ER, Bruns DE. Tietz Textbook of Clinical Chemistry and Molecular Diagnostics. 7th ed. Elsevier; 2023.
  \item Wallach J. Interpretation of Diagnostic Tests. 10th ed. Wolters Kluwer; 2022.
  \item Tierney LM, Saint S, Drazen JM. CURRENT Medical Diagnosis and Treatment. McGraw-Hill; 2023.
  \item Fischbach FT, Dunning MB. A Manual of Laboratory and Diagnostic Tests. 10th ed. Wolters Kluwer; 2023.
  \item Gennari FJ. Hypokalemia. N Engl J Med. 1998;339(7):451--458.
  \item Palmer BF, Clegg DJ. Physiology and pathophysiology of potassium homeostasis. Adv Physiol Educ. 2016;40(4):480--490.
  \item Montford JR, Linas SL. How dangerous is hyperkalemia? J Am Soc Nephrol. 2017;28(11):3155--3165.
  \item Kovesdy CP. Updates in hyperkalemia: Outcomes and therapeutic strategies. Rev Endocr Metab Disord. 2017;18(1):41--47.
\end{enumerate}

\clearpage
